\chapter{Geometry, Loading, and Design-Variable Specifications}
\label{app:geom}

This appendix lists the full numerical specification of each of
the four benchmarks used in \cref{chap:results}: nodal
coordinates, element connectivity, supports, load cases, material
constants, allowable displacements, stress limits, design-variable
groupings, and published reference optima. All values are quoted
in the benchmark's original units, which are preserved inside the
code to match the literature exactly; the 200-bar benchmark is the
only one natively specified in SI. Each table's \texttt{label}
mirrors the structure \verb|tab:geom:<bench>:<aspect>|.

\paragraph{Notation.} For every benchmark the design-variable
vector is denoted $\mathbf{a}\in\mathbb{R}^{n_{\mathrm{dv}}}$,
with each component an \emph{area group} parameter. The mapping
from group index to physical bar index is the
\emph{group-map}, stored as \texttt{group\_map} on each benchmark
class. An area-group of size $\ge 2$ ties the areas of multiple
bars together (typical in symmetric trusses). In every benchmark
the area is the only design-variable type --- sizing optimisation
only, no topology or shape modifications.

\section{10-bar Planar Cantilever}
\label{app:geom:10bar}

The 10-bar planar cantilever of
Haftka and G\"urdal~\cite{haftka_gurdal1992} is the oldest
benchmark in the sizing-optimisation literature and is
accordingly the most reported. It comprises six nodes arranged
as a $3\times 2$ rectangular grid; the two left-most nodes are
pinned supports, and two downward tip loads are applied at the
two right-most free nodes.

\begin{table}[h]
  \centering
  \small
  \caption{10-bar: nodal coordinates (inches). All nodes lie in
    the $xy$-plane.}
  \label{tab:geom:10bar:nodes}
  \begin{tabular}{rrr}
  \toprule
  Node & $X$ (in) & $Y$ (in) \\
  \midrule
  1 & 720 & 360 \\
  2 & 720 & 0 \\
  3 & 360 & 360 \\
  4 & 360 & 0 \\
  5 & 0   & 360 \\
  6 & 0   & 0   \\
  \bottomrule
  \end{tabular}
\end{table}

\begin{table}[h]
  \centering
  \small
  \caption{10-bar: element connectivity. Bar numbering follows
    the Haftka--G\"urdal convention.}
  \label{tab:geom:10bar:conn}
  \begin{tabular}{rrr}
  \toprule
  Bar & Node $i$ & Node $j$ \\
  \midrule
  1 & 5 & 3 \\ 2 & 3 & 1 \\ 3 & 6 & 4 \\ 4 & 4 & 2 \\ 5 & 3 & 4 \\
  6 & 1 & 2 \\ 7 & 4 & 5 \\ 8 & 3 & 6 \\ 9 & 2 & 3 \\ 10 & 1 & 4 \\
  \bottomrule
  \end{tabular}
\end{table}

\paragraph{Supports.} Nodes~5 and~6 are fully pinned (both $x$ and
$y$ DOFs fixed). All other nodes are free. The minimum pin count
($d(d+1)/2=3$ in 2D) is exceeded, so the structure is statically
indeterminate to first degree~\cite{subramanian2008}.

\paragraph{Loading.} A single load case is specified:
\textbf{LC1} applies a downward force of
$F_y=-\SI{100}{kip}=-\SI{100e3}{lb}$ at both node~2 and node~4.

\paragraph{Design variables.} $n_{\mathrm{dv}}=10$ --- every bar
is its own independent area variable. The area-group map is the
identity $g(i)=i$ for $i=1,\dotsc,10$.

\paragraph{Limits.} Side bounds on area:
$\num{0.1}\le a_i\le\SI{35.0}{in^2}$. Allowable displacement at
every free DOF: $|u|\le\SI{2.0}{in}$. Allowable axial stress:
$|\sigma|\le\SI{25}{ksi}=\SI{25e3}{psi}$ in both tension and
compression. Material: $E=\SI{10}{Msi}=\SI{10e6}{psi}$,
$\rho=\SI{0.1}{lb/in^3}$.

\paragraph{Reference optimum.} The canonical optimum weight
reported by Kaveh~\cite{kaveh2010} is $W^\star=\SI{5060.85}{lb}$;
our PSO reproduces this to within $\pm\SI{0.05}{lb}$ across 30
seeds (\cref{sec:results:10bar}).

\section{25-bar Spatial Transmission Tower}
\label{app:geom:25bar}

The 25-bar spatial transmission tower is a $d=3$ dimensional
benchmark with two stacked frustra. Ten nodes and twenty-five
bars are grouped into eight symmetric area groups. Two load
cases capture asymmetric and anti-symmetric wind actions.

\begin{table}[h]
  \centering
  \small
  \caption{25-bar: nodal coordinates (inches).}
  \label{tab:geom:25bar:nodes}
  \begin{tabular}{rrrr}
  \toprule
  Node & $X$ & $Y$ & $Z$ \\
  \midrule
  1 & -37.5 &  0.0 & 200 \\
  2 &  37.5 &  0.0 & 200 \\
  3 & -37.5 &  37.5 & 100 \\
  4 &  37.5 &  37.5 & 100 \\
  5 &  37.5 & -37.5 & 100 \\
  6 & -37.5 & -37.5 & 100 \\
  7 & -100  &  100  &   0 \\
  8 &  100  &  100  &   0 \\
  9 &  100  & -100  &   0 \\
  10 & -100 & -100  &   0 \\
  \bottomrule
  \end{tabular}
\end{table}

\paragraph{Supports.} Nodes 7--10 are fully pinned ($u=v=w=0$),
giving 12 prescribed DOFs and $10\cdot 3-12=18$ free DOFs.

\paragraph{Connectivity.} Bars are listed in groups reflecting
the physical role of each bar family. The 25 bars are assigned
to 8 area groups as follows:
\begin{center}
\small
\begin{tabular}{ll}
  \toprule
  Group & Bar numbers (1-indexed) \\
  \midrule
  1 (top strut)           & 1 \\
  2 (upper inclined)      & 2--5 \\
  3 (waist horizontal)    & 6--9 \\
  4 (mid vertical)        & 10--11 \\
  5 (inner diagonal)      & 12--13 \\
  6 (outer upper)         & 14--17 \\
  7 (outer middle)        & 18--21 \\
  8 (outer lower)         & 22--25 \\
  \bottomrule
\end{tabular}
\end{center}

\paragraph{Load cases.}
\begin{itemize}
  \item LC1 (asymmetric): node 1 force
    $(1000,10000,-5000)$\,lb;
    node 2 force $(0,10000,-5000)$\,lb; node 3 and node 6
    each $(500,0,0)$\,lb.
  \item LC2 (anti-symmetric): node 1 $(0,20000,-5000)$\,lb;
    node 2 $(0,-20000,-5000)$\,lb.
\end{itemize}

\paragraph{Limits.} Area bounds $\num{0.01}\le a_i\le\SI{3.4}{in^2}$.
Allowable displacement at every free DOF:
$|u|\le\SI{0.35}{in}$. Stress bounds are set per group (Table~1
of~\cite{kaveh2010}); the implementation stores them in
\texttt{group\_stress\_allow}.

\paragraph{Reference optimum.} Published Kaveh optimum
$W^\star=\SI{545.22}{lb}$; reproduced by our NSGA-II+surrogate to
within \num{0.1}\,\% across 30 seeds.

\section{72-bar Spatial Four-Tier Tower}
\label{app:geom:72bar}

The 72-bar spatial tower has twenty nodes stacked on four levels
of a $\SI{60}{in}\times\SI{60}{in}$ square footprint, each level
$\SI{60}{in}$ high. Seventy-two bars are grouped into sixteen
area groups. This is the deepest 3D benchmark in the thesis and
therefore the hardest for a constraint-obeying optimiser.

\begin{table}[h]
  \centering
  \small
  \caption{72-bar: nodal coordinates (inches). Nodes 1--16 are
    the four upper levels, nodes 17--20 the ground level.}
  \label{tab:geom:72bar:nodes}
  \begin{tabular}{rrrr}
  \toprule
  Node & $X$ & $Y$ & $Z$ \\
  \midrule
  1 & 60 & 60 & 60 \\ 2 & 60 & -60 & 60 \\
  3 & -60 & -60 & 60 \\ 4 & -60 & 60 & 60 \\
  5 & 60 & 60 & 120 \\ 6 & 60 & -60 & 120 \\
  7 & -60 & -60 & 120 \\ 8 & -60 & 60 & 120 \\
  9 & 60 & 60 & 180 \\ 10 & 60 & -60 & 180 \\
  11 & -60 & -60 & 180 \\ 12 & -60 & 60 & 180 \\
  13 & 60 & 60 & 240 \\ 14 & 60 & -60 & 240 \\
  15 & -60 & -60 & 240 \\ 16 & -60 & 60 & 240 \\
  17 & 60 & 60 & 0 \\ 18 & 60 & -60 & 0 \\
  19 & -60 & -60 & 0 \\ 20 & -60 & 60 & 0 \\
  \bottomrule
  \end{tabular}
\end{table}

\begin{table}[h]
  \centering
  \small
  \caption{72-bar: element connectivity (first 20 of 72). The
    remaining 52 bars follow the four-fold tier-symmetric pattern
    of the first tier --- see \texttt{src/benchmarks/truss\_72bar.py}
    for the complete table.}
  \label{tab:geom:72bar:conn}
  \begin{tabular}{rrr}
  \toprule
  Bar & Node $i$ & Node $j$ \\
  \midrule
  1 & 17 & 1 \\ 2 & 18 & 2 \\ 3 & 19 & 3 \\ 4 & 20 & 4 \\
  5 & 1 & 2 \\ 6 & 2 & 3 \\ 7 & 3 & 4 \\ 8 & 4 & 1 \\
  9 & 1 & 18 \\ 10 & 2 & 19 \\ 11 & 3 & 20 \\ 12 & 4 & 17 \\
  13 & 17 & 2 \\ 14 & 18 & 3 \\ 15 & 19 & 4 \\ 16 & 20 & 1 \\
  17 & 1 & 3 \\ 18 & 2 & 4 \\ 19 & 1 & 5 \\ 20 & 2 & 6 \\
  \multicolumn{3}{c}{$\cdots$ 52 rows omitted, see source} \\
  \bottomrule
  \end{tabular}
\end{table}

\paragraph{Supports.} Nodes 17--20 are fully pinned. The other
16 nodes are free, giving 48 free DOFs.

\paragraph{Load cases.}
\begin{itemize}
  \item LC1: node~13 force $(5000,5000,-5000)$\,lb
    (asymmetric tip load at one corner).
  \item LC2: $-\SI{5000}{lb}$ vertical at each of the four
    top-level nodes 13--16.
\end{itemize}

\paragraph{Design variables and groups.} $n_{\mathrm{dv}}=16$.
Each of the four tiers has four area-groups (bottom columns,
horizontals, verticals, and diagonals), giving $4\times 4=16$.

\paragraph{Limits.} Area bounds $\num{0.1}\le a_i\le\SI{4.0}{in^2}$.
Allowable displacement at every top-level free DOF
$(u_x, u_y$ at nodes 13--15$)$: $|u|\le\SI{0.25}{in}$.
Stress limit as 10-bar.

\paragraph{Reference optimum.}
$W^\star=\SI{379.62}{lb}$ (Kaveh and Talatahari~\cite{kaveh2010}).

\section{200-bar Planar Stepped Tower}
\label{app:geom:200bar}

The 200-bar benchmark is the largest truss in the thesis and the
only benchmark defined in SI units. It is a 2D stepped
seven-bay-wide tower of eleven full-width tiers plus the top
apex; its geometry is reproduced from the original
Kaveh--Kalateh-Ahani formulation~\cite{kaveh2011}, but
coordinates are our own cm-rounded coordinates (indicated in the
benchmark API with \texttt{reference\_verified=False}).

\begin{table}[h]
  \centering
  \small
  \caption{200-bar: nodal coordinates (metres). First 15 of 77
    total nodes shown; remaining nodes continue the vertical
    $\SI{3}{m}$ stride pattern.}
  \label{tab:geom:200bar:nodes}
  \begin{tabular}{rrr}
  \toprule
  Node & $X$ (m) & $Y$ (m) \\
  \midrule
  1 & -4.5 & 0 \\ 2 & -3.0 & 0 \\ 3 & -1.5 & 0 \\
  4 & 0 & 0 \\ 5 & 1.5 & 0 \\ 6 & 3.0 & 0 \\ 7 & 4.5 & 0 \\
  8 & -4.5 & 3 \\ 9 & -3.0 & 3 \\ 10 & -1.5 & 3 \\
  11 & 0 & 3 \\ 12 & 1.5 & 3 \\ 13 & 3.0 & 3 \\ 14 & 4.5 & 3 \\
  15 & -4.5 & 6 \\
  \multicolumn{3}{c}{$\cdots$ 62 rows omitted, see source} \\
  \bottomrule
  \end{tabular}
\end{table}

\paragraph{Supports.} Nodes 1--7 (bottom row) are fully pinned in
both translations: 14 DOFs fixed out of $77\times 2=154$ total,
leaving 140 free.

\paragraph{Load cases.} Three gravity load cases are specified:
\begin{itemize}
  \item LC1: $\SI{1}{kN}$ downward at every top-level node
    (nodes 71--77).
  \item LC2: $\SI{10}{kN}$ downward at top-level nodes 71--73
    and at every east-side edge node (nodes 8, 22, 36, 50, 64).
  \item LC3: LC1 + LC2 superposed.
\end{itemize}

\paragraph{Design variables.} $n_{\mathrm{dv}}=29$. Each group
represents one structural role (edge chord, internal vertical,
diagonal, horizontal) at one tier of the tower; the full
group-map is listed in \texttt{src/benchmarks/truss\_200bar.py}.

\paragraph{Limits.} Area bounds $\num{6.5e-5}\le a_i\le
\SI{2.5e-2}{m^2}$ (equivalent to
$\SI{0.1}\le a_i\le\SI{38.75}{in^2}$ --- wide enough to admit
the same search space as the original imperial formulation).
Displacement limit: every top-row free DOF is capped at
$\SI{0.1}{m}$. Material: $E=\SI{210}{GPa}$,
$\rho=\SI{7850}{kg/m^3}$ (structural steel). Stress allowable:
$|\sigma|\le\SI{250}{MPa}$ per Fe~410.

\paragraph{Reference optimum.} No published SI-unit optimum
matches our coordinate choices. Our own batch of 30
PSO seeds (Ch.~\ref{chap:results} Section~\ref{sec:results:200bar})
converges to $W^\star\approx\SI{315.89}{kg}$ for the best seed,
with median $\SI{330.1}{kg}$. These serve as the reference for
Phase 10 ablation studies and are reported to four significant
figures.

\section{Cross-Benchmark Summary}
\label{app:geom:summary}

The four benchmarks span two orders of magnitude in number of
design variables and bars, and both 2D and 3D configurations,
giving a balanced coverage of the sizing-optimisation difficulty
landscape (\cref{tab:geom:summary}).

\begin{table}[h]
  \centering
  \caption{Cross-benchmark numerical summary. Units kept in each
    benchmark's native system for clarity.}
  \label{tab:geom:summary}
  \small
  \begin{tabular}{lcccc}
    \toprule
    Quantity & 10-bar & 25-bar & 72-bar & 200-bar \\
    \midrule
    Dimension $d$          & 2     & 3     & 3     & 2 \\
    Nodes $N$              & 6     & 10    & 20    & 77 \\
    Bars $M$               & 10    & 25    & 72    & 200 \\
    Design vars $n_\mathrm{dv}$ & 10  & 8   & 16   & 29 \\
    Load cases             & 1     & 2     & 2     & 3 \\
    Displacement limit     & $\SI{2.0}{in}$ & $\SI{0.35}{in}$ & $\SI{0.25}{in}$ & $\SI{0.1}{m}$ \\
    Stress allowable       & $\SI{25}{ksi}$ & (grouped) & $\SI{25}{ksi}$ & $\SI{250}{MPa}$ \\
    Units                  & Imperial & Imperial & Imperial & SI \\
    Reference $W^\star$     & $\SI{5060.85}{lb}$ & $\SI{545.22}{lb}$ & $\SI{379.62}{lb}$ & $\SI{315.89}{kg}$ \\
    Pinned DOFs            & 4     & 12    & 12    & 14 \\
    Free DOFs              & 8     & 18    & 48    & 140 \\
    \bottomrule
  \end{tabular}
\end{table}

\paragraph{Closing remark.} Although the problem sizes differ by
over an order of magnitude, the same constraint API
(\texttt{src/constraints/}) and the same FEM backend
(\cref{app:fem}) handle all four without benchmark-specific
logic. The only benchmark-specific information is the contents of
\texttt{src/benchmarks/truss\_*.py} files; those files
collectively reproduce the tables above and supply them to the
generic evaluator via \texttt{BenchmarkBase.evaluate()}.
